\chapter{Range-Extended Network-Aware Model Partitioning Framework} \label{chap:ch4}

\section{Approach} \label{sec:se41}

After analyzing the state of the art on model partitioning and \ac{ML} 
orchestration in edge computing, we found that current frameworks tend to 
make use of less constrained devices and use more powerful and even less 
constrained devices only in the later stages of the inference tasks, which may 
not be the most optimal approach in a network when the goal is to make full use 
of all available devices, similar to a private local cluster. Additionally, there is a 
lack of available partitioning methods tackling models with cyclic topology 
(e.g., \acp{RNN}, MobileNets, etc.) to the edge environment. 
% TODO add FR's examples

To address these issues, we propose a novel model partitioning framework for 
efficient integration of edge devices within far-edge device networks, enabling 
collaborative inference in highly constrained environments. The framework primarily 
leverages the densely interconnected deployed far-edge devices, with minimal 
reliance on more capable edge devices. Unlike typical edge deployments that 
distribute work across heterogeneous nodes separated by non-negligible 
communication delays, this approach targets networks where devices experience 
negligible communication delays, creating a more challenging but realistic 
scenario for far-edge computing.
Additionally, the framework can exploit the cyclic topology of certain \ac{ML} 
models (e.g., \acp{RNN}), enabling device reuse across multiple inference stages. 
By leveraging the recurrent nature of these models, individual devices can process 
different layers or iterations, effectively acting as shared computational resources 
throughout the inference pipeline.

\section{Methods} \label{sec:se42}

For the testing and validation stages of the proposed framework, we will use a set 
of 15 \ac{MCU}-based devices with available memory on the order of kilobytes and 
one \ac{MPU}-based device with available memory on the order of megabytes. 

The devices will be connected via Wi-Fi and deployed in an 
indoor environment. The topology represents a typical smart building 
scenario, forming a fully-connected graph of 16 nodes.

The metrics for evaluation will include inference latency, energy consumption, and
memory usage per device. We will then compare the performance of our proposed
framework against the state of the art model partitioning frameworks that we
analysed in Chapter~\ref{chap:ch3}.

\section{Expected Results} \label{sec:se43}

We expect to achieve a successful deployment of partitioned models with cyclic 
components, extending the range of \ac{DNN} models that can be deployed in edge
environments. 

Additionally, we anticipate improved performance in terms of
latency and energy efficiency compared to existing model partitioning frameworks,
especially in scenarios with a high density of constrained devices. This is
partially due to improved communication methods between devices and more efficient
utilization of the available computational resources across the network, in
particular, the use of the edge device among the far-edge devices.

\section{Project plan} \label{sec:se44}

The project is planned to be executed over a period of 5 months, resulting in 
18 available weeks, divided into the following tasks:

\begin{description}[style=unboxed, leftmargin=1cm, font=\bfseries]

    \item[Testbed Setup (Weeks 1-8):]
    This task involves setting up the hardware and software environment for testing,
    including the configuration of the devices and network. Furthermore, we will
    implement the necessary tools for monitoring and evaluating performance metrics.
    
    \item[Deployment of the Model Partition Segments (Weeks 4-7):]
    This task involves deploying the partitioned model segments onto the
    selected devices, ensuring correct functionality and efficient communication
    between them.
    
    \item[Partitioning Subprocedure Integration (Weeks 7-12):]
    This task focuses on integrating a core partitioning algorithm, chosen from 
    an existing framework reviewed in chapter~\ref{chap:ch3}, that will divide
    the \ac{DNN} models into lightweight segments suitable for the constrained 
    devices in the network. Ideally, this subprocedure will be extended to take 
    into account the cyclic topology of certain models to optimize the partitioning 
    process.
    
    \item[Testing and Evaluation (Weeks 12-14):]
    During these weeks, we will conduct extensive testing of the deployed framework,
    measuring the defined performance metrics and comparing them against baseline
    results from existing frameworks.

    \item[Dissertation Elaboration (Weeks 15-18):]
    After completing the implementation work and testing, we will focus on writing
    the dissertation, documenting the methodology, results, and conclusions of the
    project.

    
\end{description}

The first three tasks represent the milestone corresponding to the MVP of the
project, expected to be completed by the end of week 12. The testing and evaluation
task will provide the necessary data to analyze the performance of the proposed
framework, while the dissertation elaboration task will ensure that the findings
are well-documented and presented. The workflow can be better visualized in the Gantt
chart shown in Figure~\ref{fig:gantt}. It should be noted that the weeks corresponding
to the Easter break and Academic Week are not taken into account in this timeline.

\begin{figure}[H]
  \centering
  \includegraphics[width=1\textwidth]{approach-gantt.png}
  \caption{Gantt chart representing the project plan and timeline.}
  \label{fig:gantt}
\end{figure}

\section{Conclusions} \label{sec:se45}

This chapter proposes a novel framework to efficiently execute collaborative 
inference tasks,
excelling in the integration of edge devices within a far-edge device network and 
managing to make use of highly constrained devices. Additionally, the framework 
extends support to \ac{DNN} models with cyclic topology through device reuse across 
inference stages. This marks a significant advancement in the field of edge computing
and \ac{IoT}, opening new possibilities for applications that require real-time 
processing and low latency.
